\documentclass[12pt]{article}
\usepackage[utf8]{inputenc}
\usepackage[T1]{fontenc}
\usepackage{amsmath,amssymb}
\usepackage{geometry}
\geometry{margin=2.5cm}

\begin{document}

\noindent $\tau \overset{\text{def}}{=} (i,j)$

\noindent $\pi \overset{\text{def}}{=} \sigma\circ\tau$

\[
\pi(j) = \sigma\big(\tau(j)\big) = \sigma(i) = j
\]

\bigskip
\noindent Dacă $\sigma(k)=k \Rightarrow \pi(k)=k \Rightarrow k\neq i, k\neq j \ \Rightarrow d_\pi<d_\sigma.$

\[
\pi(k) = \sigma\big(\tau(k)\big) = \sigma(k) = k.
\]

\noindent $\pi = \sigma\circ\tau \ \big| \ \tau.$

\bigskip
\noindent $\pi = \tau_1\circ\cdots\circ\tau_s$ \ (ip.\ de ind.) \ $\big|\tau \ \Rightarrow \ \tau_1\circ\cdots\circ\tau_s\circ\tau=\sigma.$

\bigskip
\noindent Dacă $\sigma\in S_n$ \ inv$(\sigma) \overset{\text{not}}{=} \text{card}\{(i,j) \mid 1\leq i<j\leq n,\ \sigma(i)>\sigma(j)\}.$

\noindent (= nr.\ inversiunilor lui $\sigma$.)

\[
\sigma = \begin{pmatrix}1&2&\cdots&i&=&j&\cdots&n\\ \sigma(1)&\sigma(2)&\cdots&\sigma(i)&\cdots&\sigma(j)&\cdots&\sigma(n)\end{pmatrix} \quad C_n^2.
\]

\noindent inv$(e)=0$.

\bigskip
\noindent\fbox{Def}: \ Nr.\ $\varepsilon(\sigma) = (-1)^{\text{inv}(\sigma)}$ s.n.\ \underline{signatura}.

\bigskip
\noindent $\sigma$=pară dacă $\varepsilon(\sigma)=1$.

\noindent \phantom{$\sigma$=}impară dacă $\varepsilon(\sigma)=-1$.

\bigskip
\noindent\fbox{T} \ Dacă $\sigma\in S_n$ şi $\sigma=\tau_1\circ\tau_2\circ\cdots\circ\tau_m \Rightarrow$ nr.\ $m$ şi inv$(\sigma)$ au aceeaşi

\noindent paritate.

\bigskip
\noindent\fbox{Lemă} \ (a) \ Fie $\tau,\sigma\in S_n$, $\tau$=transpoziţie

\noindent Atunci inv$(\sigma)$ şi inv$(\sigma\circ\tau)$ au parităţi diferite.

\noindent (b) $(\forall)$ transp.\ e permut.\ impară.

\bigskip
\noindent\fbox{Dem\ Lemă} \ $\tau=(i,j)$ \quad $1\leq i<j\leq n$ şi $r=j-i\in\mathbb{N}^*$.

\bigskip
\noindent Dacă $r=1$

\[
\sigma = \begin{pmatrix}1&\cdots&i&i+1&\cdots&n\\ \sigma(1)&\cdots&\sigma(i)&\sigma(i+1)&\cdots&\sigma(n)\end{pmatrix}
= \begin{pmatrix}1&\cdots&i&j&\cdots&n\\ \sigma(1)&\cdots&\sigma(i)&\sigma(j)&\cdots&\sigma(n)\end{pmatrix}
\]

\[
\pi=\sigma\circ\tau = \begin{pmatrix}1&\cdots&i&j&\cdots&n\\ \sigma(1)&\cdots&\sigma(j)&\sigma(i)&\cdots&\sigma(n)\end{pmatrix}
= \begin{pmatrix}1&\cdots&i&i+1&\cdots&n\\ \sigma(1)&\cdots&\sigma(i+1)&\sigma(i)&\cdots&\sigma(n)\end{pmatrix}
\]

\bigskip
\noindent $r=1 \Rightarrow$ inv$(\sigma\circ\tau) =
\begin{cases}
\text{inv}(\sigma)+1, & \text{dacă } \sigma(i)<\sigma(i+1)\\
\text{inv}(\sigma)-1, & \text{dacă } \sigma(i)>\sigma(i+1)
\end{cases}$

\bigskip
\noindent\underline{Cazul general}: \ $\tau=(i,j) \Rightarrow j-i=r$, $r>1$.

\noindent $\pi$ se obţine din $\sigma$ prin schimbarea lui $\sigma(i)$ cu $\sigma(j)$, realizată printr-o succesiune de transpoziţii de poziţii adiacente.\footnote{acest paragraf al demonstraţiei e foarte compactat în original; pasul exact (,,$d_\pi$ ... schimb ...'') n-a putut fi reconstituit cu certitudine deplină, dar argumentul standard e cel prezentat aici.}

\[
2r-1 \text{ etape}
\]

\noindent fiecare etapă schimbă paritatea lui inv (cazul $r=1$) $\Rightarrow$ după $2r-1$ etape (nr.\ impar)

\noindent paritatea lui inv$(\sigma)$ se schimbă. \hfill parit.\ diferite \quad Q.E.D.

\end{document}
